\chapter{Appendix to Chapter 3: Detailed Proofs}
\label{app:ch3_proofs}

This appendix provides the algebraic details underlying the main results in
Chapter~\ref{chap:model1}. To keep the main chapter focused on economic
interpretation, the text there reports the closed-form equilibria and the main
comparative insights. Here we show the intermediate steps used to obtain those
results.

\section{Useful Re-parameterization and Preliminary Observations}
\label{app:preliminaries}

For the decentralized games it is convenient to work with the retailer markup $m = p-w,$ so that the retail price can be written as $ p=w+m.$
Substituting into the demand system $Q(p,s)$ from given in  \ref{eq:demand} 
% \begin{equation}
%  Q(p,s)=\bar Q-\alpha p+\beta s,
% \end{equation}
yields  $ Q(w,m,s)=\bar Q-\alpha(w+m)+\beta s.$
% \begin{equation*}
%  Q(w,m,s)=\bar Q-\alpha(w+m)+\beta s.
%  \label{eq:app_demand_markup}
% \end{equation*}

The manufacturer profit becomes
\begin{equation}
 \Pi_M(w,m,s)=w\bigl(\bar Q-\alpha(w+m)+\beta s\bigr),
 \label{eq:app_piM}
\end{equation}
while the retailer profit becomes
\begin{equation}
 \Pi_R(w,m,s)=m\bigl(\bar Q-\alpha(w+m)+\beta s\bigr)-\frac{1}{2}\gamma s^2.
 \label{eq:app_piR}
\end{equation}

The quadratic structure makes the second-order conditions transparent, but the
conditions are regime-specific. In VN and in the retailer follower problem under
MS, the retailer's Hessian in $(m,s)$ is
\begin{equation}
 H_R^{VN/MS}=
 \begin{pmatrix}
 -2\alpha & \beta\\
 \beta & -\gamma
 \end{pmatrix},
\end{equation}
which is negative definite if and only if
\begin{equation}
 \Delta_2:=2\alpha\gamma-\beta^2>0.
\end{equation}
The integrated benchmark has the same Hessian form in $(p,s)$ and therefore
uses the same condition. By contrast, after the manufacturer's response is
substituted into the RS game, the retailer leader's reduced problem has Hessian
\begin{equation}
 H_R^{RS}=
 \begin{pmatrix}
 -\alpha & \beta/2\\
 \beta/2 & -\gamma
 \end{pmatrix},
\end{equation}
which is negative definite if and only if
\begin{equation}
 \Delta_4:=4\alpha\gamma-\beta^2>0.
\end{equation}
Thus denominator positivity and second-order sufficiency are distinct. For
example, $\Delta_3>0$ makes the VN closed-form denominator positive, but the
retailer problem still requires $\Delta_2>0$ for the stationary point to be a
maximum. The cross-regime comparisons in Chapter~\ref{chap:model1} therefore
maintain $\Delta_2>0$, while RS is also valid as a standalone game on the wider
domain $\Delta_4>0$.

\section{Proof of Proposition~\ref{prop:vn}: Vertical Nash Equilibrium}
\label{app:proof_vn}

Under Vertical Nash, the manufacturer chooses $w$ and the retailer chooses
$(m,s)$ simultaneously. Using \eqref{eq:app_piM}--\eqref{eq:app_piR}, the
players solve
\begin{equation}
 \max_w \Pi_M(w,m,s),
 \qquad
 \max_{m,s}\Pi_R(w,m,s).
\end{equation}

\subsection{Step 1: Retailer first-order conditions}

Differentiating \eqref{eq:app_piR} with respect to $m$ and $s$ gives
\begin{align}
 \frac{\partial \Pi_R}{\partial m}
 &= \bar Q-\alpha(w+2m)+\beta s = 0,
 \label{eq:app_vn_foc_m}\\
 \frac{\partial \Pi_R}{\partial s}
 &= \beta m-\gamma s = 0.
 \label{eq:app_vn_foc_s}
\end{align}
Equation~\eqref{eq:app_vn_foc_s} implies the retailer's recycling maturity is
proportional to its markup:
\begin{equation}
 s=\frac{\beta}{\gamma}m.
 \label{eq:app_vn_s_of_m}
\end{equation}

\subsection{Step 2: Manufacturer first-order condition}

Differentiating \eqref{eq:app_piM} with respect to $w$ yields
\begin{equation}
 \frac{\partial \Pi_M}{\partial w}
 = \bar Q-\alpha(2w+m)+\beta s = 0.
 \label{eq:app_vn_foc_w}
\end{equation}

\subsection{Step 3: Solve the system}

Subtract \eqref{eq:app_vn_foc_w} from \eqref{eq:app_vn_foc_m}:
\begin{equation}
\bigl[\bar Q-\alpha(w+2m)+\beta s\bigr]-\bigl[\bar Q-\alpha(2w+m)+\beta s\bigr]=0,
\end{equation}
which simplifies to
\begin{equation}
 \alpha(w-m)=0.
\end{equation}
Hence
\begin{equation}
 w^{VN}=m^{VN}.
 \label{eq:app_vn_equal_margins}
\end{equation}

Now substitute \eqref{eq:app_vn_equal_margins} and
\eqref{eq:app_vn_s_of_m} into \eqref{eq:app_vn_foc_w}:
\begin{equation}
 \bar Q-\alpha(3w)+\beta\left(\frac{\beta}{\gamma}w\right)=0.
\end{equation}
Multiplying by $\gamma$ gives
\begin{equation}
 \bar Q\gamma-(3\alpha\gamma-\beta^2)w=0,
\end{equation}
so
\begin{equation}
 w^{VN}=\frac{\bar Q\gamma}{3\alpha\gamma-\beta^2}.
\end{equation}
From \eqref{eq:app_vn_equal_margins},
\begin{equation}
 m^{VN}=\frac{\bar Q\gamma}{3\alpha\gamma-\beta^2}.
\end{equation}
Using \eqref{eq:app_vn_s_of_m},
\begin{equation}
 s^{VN}=\frac{\bar Q\beta}{3\alpha\gamma-\beta^2}.
\end{equation}
Since $p=w+m$,
\begin{equation}
 p^{VN}=\frac{2\bar Q\gamma}{3\alpha\gamma-\beta^2}.
\end{equation}

\subsection{Step 4: Demand and profits}

Substitute the equilibrium values into demand:
\begin{align}
 Q^{VN}
 &= \bar Q-\alpha p^{VN}+\beta s^{VN} \\
 &= \bar Q-\alpha\left(\frac{2\bar Q\gamma}{3\alpha\gamma-\beta^2}\right)
 +\beta\left(\frac{\bar Q\beta}{3\alpha\gamma-\beta^2}\right) \\
 &= \frac{\alpha\gamma\bar Q}{3\alpha\gamma-\beta^2}.
\end{align}
Therefore,
\begin{equation}
 \Pi_M^{VN}=w^{VN}Q^{VN}
 =\frac{\alpha\bar Q^2\gamma^2}{(3\alpha\gamma-\beta^2)^2},
\end{equation}
and
\begin{align}
 \Pi_R^{VN}
 &=m^{VN}Q^{VN}-\frac{1}{2}\gamma(s^{VN})^2 \\
 &=\frac{\alpha\bar Q^2\gamma^2}{(3\alpha\gamma-\beta^2)^2}
 -\frac{1}{2}\gamma\left(\frac{\bar Q\beta}{3\alpha\gamma-\beta^2}\right)^2 \\
 &=\frac{\gamma\bar Q^2(2\alpha\gamma-\beta^2)}{2(3\alpha\gamma-\beta^2)^2}.
\end{align}
This proves Proposition~\ref{prop:vn}.

\section{Proof of Proposition~\ref{prop:ms}: Manufacturer Stackelberg Equilibrium}
\label{app:proof_ms}

Under MS, the manufacturer chooses $w$ first, and the retailer then solves for
its best response in $(m,s)$.

\subsection{Step 1: Retailer best response given $w$}

Given $w$, the retailer solves
\begin{equation}
 \max_{m,s}
 \left\{m\bigl(\bar Q-\alpha(w+m)+\beta s\bigr)-\frac{1}{2}\gamma s^2\right\}.
\end{equation}
The first-order conditions are
\begin{align}
 \bar Q-\alpha w-2\alpha m+\beta s &=0,
 \label{eq:app_ms_foc_m}\\
 \beta m-\gamma s &=0.
 \label{eq:app_ms_foc_s}
\end{align}
From \eqref{eq:app_ms_foc_s},
\begin{equation}
 s=\frac{\beta}{\gamma}m.
 \label{eq:app_ms_s_of_m}
\end{equation}
Substituting into \eqref{eq:app_ms_foc_m} gives
\begin{equation}
 \bar Q-\alpha w-2\alpha m+\frac{\beta^2}{\gamma}m=0.
\end{equation}
Multiplying by $\gamma$,
\begin{equation}
 \gamma(\bar Q-\alpha w)-(2\alpha\gamma-\beta^2)m=0.
\end{equation}
Hence the retailer's best-response functions are
\begin{equation}
 m^{BR}(w)=\frac{\gamma(\bar Q-\alpha w)}{2\alpha\gamma-\beta^2},
 \qquad
 s^{BR}(w)=\frac{\beta(\bar Q-\alpha w)}{2\alpha\gamma-\beta^2}.
 \label{eq:app_ms_br}
\end{equation}

\subsection{Step 2: Manufacturer problem}

Substitute \eqref{eq:app_ms_br} into the demand expression:
\begin{align}
 Q^{BR}(w)
 &= \bar Q-\alpha\bigl(w+m^{BR}(w)\bigr)+\beta s^{BR}(w) \\
 &= \bar Q-\alpha w
 -\alpha\frac{\gamma(\bar Q-\alpha w)}{2\alpha\gamma-\beta^2}
 +\beta\frac{\beta(\bar Q-\alpha w)}{2\alpha\gamma-\beta^2} \\
 &= \frac{\alpha\gamma(\bar Q-\alpha w)}{2\alpha\gamma-\beta^2}.
\end{align}
Therefore the manufacturer solves
\begin{equation}
 \max_w \Pi_M(w)=w\cdot \frac{\alpha\gamma(\bar Q-\alpha w)}{2\alpha\gamma-\beta^2}.
\end{equation}
Because the denominator is constant in $w$, this is equivalent to maximizing
\begin{equation}
 f(w)=w(\bar Q-\alpha w)=\bar Q w-\alpha w^2.
\end{equation}
The derivative is
\begin{equation}
 f'(w)=\bar Q-2\alpha w.
\end{equation}
Setting $f'(w)=0$ gives
\begin{equation}
 w^{MS}=\frac{\bar Q}{2\alpha}.
\end{equation}
Since $f''(w)=-2\alpha<0$, this is the unique optimum.

\subsection{Step 3: Follower decisions}

Substitute $w^{MS}$ into \eqref{eq:app_ms_br}:
\begin{align}
 m^{MS}
 &=\frac{\gamma\left(\bar Q-\alpha\cdot\frac{\bar Q}{2\alpha}\right)}{2\alpha\gamma-\beta^2}
 =\frac{\bar Q\gamma}{2(2\alpha\gamma-\beta^2)},\\
 s^{MS}
 &=\frac{\beta\left(\bar Q-\alpha\cdot\frac{\bar Q}{2\alpha}\right)}{2\alpha\gamma-\beta^2}
 =\frac{\bar Q\beta}{2(2\alpha\gamma-\beta^2)}.
\end{align}
Then
\begin{equation}
 p^{MS}=w^{MS}+m^{MS}
 =\frac{\bar Q}{2\alpha}+\frac{\bar Q\gamma}{2(2\alpha\gamma-\beta^2)}
 =\frac{\bar Q(3\alpha\gamma-\beta^2)}{2\alpha(2\alpha\gamma-\beta^2)}.
\end{equation}

\subsection{Step 4: Demand and profits}

Demand is
\begin{equation}
 Q^{MS}=\frac{\alpha\gamma\bar Q}{2(2\alpha\gamma-\beta^2)}.
\end{equation}
Hence
\begin{equation}
 \Pi_M^{MS}=w^{MS}Q^{MS}
 =\frac{\bar Q}{2\alpha}\cdot \frac{\alpha\gamma\bar Q}{2(2\alpha\gamma-\beta^2)}
 =\frac{\gamma\bar Q^2}{4(2\alpha\gamma-\beta^2)},
\end{equation}
and
\begin{equation}
 \Pi_R^{MS}=m^{MS}Q^{MS}-\frac{1}{2}\gamma(s^{MS})^2
 =\frac{\gamma\bar Q^2}{8(2\alpha\gamma-\beta^2)}.
\end{equation}
This proves Proposition~\ref{prop:ms}.

\section{Proof of Proposition~\ref{prop:rs}: Retailer Stackelberg Equilibrium}
\label{app:proof_rs}

Under RS, the retailer chooses $(m,s)$ first and the manufacturer responds in
$w$.

\subsection{Step 1: Manufacturer best response given $(m,s)$}

Given $(m,s)$, the manufacturer solves
\begin{equation}
 \max_w\; w\bigl(\bar Q-\alpha(w+m)+\beta s\bigr).
\end{equation}
The first-order condition is
\begin{equation}
 \bar Q-2\alpha w-\alpha m+\beta s=0,
\end{equation}
which yields the best response
\begin{equation}
 w^{BR}(m,s)=\frac{\bar Q-\alpha m+\beta s}{2\alpha}.
 \label{eq:app_rs_brw}
\end{equation}
Because the second derivative is $-2\alpha<0$, this is optimal.

\subsection{Step 2: Retailer reduced-form problem}

Substitute \eqref{eq:app_rs_brw} into demand:
\begin{align}
 Q &= \bar Q-\alpha\bigl(w^{BR}(m,s)+m\bigr)+\beta s \\
   &= \bar Q-\alpha\left(\frac{\bar Q-\alpha m+\beta s}{2\alpha}+m\right)+\beta s \\
   &= \frac{\bar Q-\alpha m+\beta s}{2}.
 \label{eq:app_rs_reduced_demand}
\end{align}
Thus the retailer's reduced-form objective is
\begin{align}
 \Pi_R(m,s)
 &=m\cdot \frac{\bar Q-\alpha m+\beta s}{2}-\frac{1}{2}\gamma s^2.
 \label{eq:app_rs_reduced_obj}
\end{align}
The Hessian of this reduced-form objective is
\begin{equation}
 H_R^{RS}=
 \begin{pmatrix}
 -\alpha & \beta/2\\
 \beta/2 & -\gamma
 \end{pmatrix}.
\end{equation}
It is negative definite if and only if $4\alpha\gamma-\beta^2>0$. Under this
condition, the first-order conditions are sufficient for the unique interior RS
equilibrium.
The first-order conditions are
\begin{align}
 \frac{\partial \Pi_R}{\partial m}
 &=\frac{\bar Q-2\alpha m+\beta s}{2}=0,
 \label{eq:app_rs_foc_m}\\
 \frac{\partial \Pi_R}{\partial s}
 &=\frac{\beta m}{2}-\gamma s=0.
 \label{eq:app_rs_foc_s}
\end{align}
From \eqref{eq:app_rs_foc_s},
\begin{equation}
 s=\frac{\beta m}{2\gamma}.
 \label{eq:app_rs_s_of_m}
\end{equation}
Substituting into \eqref{eq:app_rs_foc_m} gives
\begin{equation}
 \bar Q-2\alpha m+\frac{\beta^2}{2\gamma}m=0.
\end{equation}
Multiplying by $2\gamma$,
\begin{equation}
 2\gamma\bar Q-(4\alpha\gamma-\beta^2)m=0.
\end{equation}
Hence
\begin{equation}
 m^{RS}=\frac{2\gamma\bar Q}{4\alpha\gamma-\beta^2},
 \qquad
 s^{RS}=\frac{\beta\bar Q}{4\alpha\gamma-\beta^2}.
\end{equation}

\subsection{Step 3: Manufacturer response, demand, and profits}

Substituting the retailer's choices into \eqref{eq:app_rs_brw},
\begin{align}
 w^{RS}
 &=\frac{\bar Q-\alpha\cdot \frac{2\gamma\bar Q}{4\alpha\gamma-\beta^2}
 +\beta\cdot\frac{\beta\bar Q}{4\alpha\gamma-\beta^2}}{2\alpha} \\
 &=\frac{\gamma\bar Q}{4\alpha\gamma-\beta^2}.
\end{align}
Therefore
\begin{equation}
 p^{RS}=w^{RS}+m^{RS}=\frac{3\gamma\bar Q}{4\alpha\gamma-\beta^2}.
\end{equation}
Demand is
\begin{equation}
 Q^{RS}=\bar Q-\alpha p^{RS}+\beta s^{RS}
 =\frac{\alpha\gamma\bar Q}{4\alpha\gamma-\beta^2}.
\end{equation}
Thus
\begin{equation}
 \Pi_M^{RS}=w^{RS}Q^{RS}
 =\frac{\alpha\gamma^2\bar Q^2}{(4\alpha\gamma-\beta^2)^2},
\end{equation}
and
\begin{equation}
 \Pi_R^{RS}=m^{RS}Q^{RS}-\frac{1}{2}\gamma(s^{RS})^2
 =\frac{\gamma\bar Q^2}{2(4\alpha\gamma-\beta^2)}.
\end{equation}
This proves Proposition~\ref{prop:rs}.

\section{Proof of Proposition~\ref{prop:in}: Integrated Benchmark}
\label{app:proof_in}

Under integration, the channel chooses $(p,s)$ to maximize total profit
\begin{equation}
 \Pi_{Total}(p,s)=p(\bar Q-\alpha p+\beta s)-\frac{1}{2}\gamma s^2.
\end{equation}
The first-order conditions are
\begin{align}
 \bar Q-2\alpha p+\beta s&=0,
 \label{eq:app_in_foc_p}\\
 \beta p-\gamma s&=0.
 \label{eq:app_in_foc_s}
\end{align}
From \eqref{eq:app_in_foc_s},
\begin{equation}
 s=\frac{\beta}{\gamma}p.
\end{equation}
Substitute into \eqref{eq:app_in_foc_p}:
\begin{equation}
 \bar Q-2\alpha p+\frac{\beta^2}{\gamma}p=0.
\end{equation}
Multiplying by $\gamma$,
\begin{equation}
 \bar Q\gamma-(2\alpha\gamma-\beta^2)p=0,
\end{equation}
so
\begin{equation}
 p^{IN}=\frac{\bar Q\gamma}{2\alpha\gamma-\beta^2},
 \qquad
 s^{IN}=\frac{\bar Q\beta}{2\alpha\gamma-\beta^2}.
\end{equation}
Demand is
\begin{equation}
 Q^{IN}=\bar Q-\alpha p^{IN}+\beta s^{IN}
 =\frac{\alpha\gamma\bar Q}{2\alpha\gamma-\beta^2}.
\end{equation}
Finally,
\begin{equation}
 \Pi_{Total}^{IN}=p^{IN}Q^{IN}-\frac{1}{2}\gamma(s^{IN})^2
 =\frac{\gamma\bar Q^2}{2(2\alpha\gamma-\beta^2)}.
\end{equation}
Under the equal-split normalization used in Chapter~\ref{chap:model1},
\begin{equation}
 \Pi_M^{IN}=\Pi_R^{IN}=\frac{1}{2}\Pi_{Total}^{IN}
 =\frac{\gamma\bar Q^2}{4(2\alpha\gamma-\beta^2)}.
\end{equation}
This proves Proposition~\ref{prop:in}.

\section{Auxiliary Derivations for Cross-Regime Comparisons}
\label{app:rankings}

This section records the algebra behind the ranking results summarized in
Chapter~\ref{chap:model1}. Let
\begin{equation}
 x:=\frac{\beta^2}{\alpha\gamma},
 \qquad 0<x<2,
\end{equation}
so that interior solutions for all regimes are available. Then
\begin{equation}
 \Delta_2=\alpha\gamma(2-x),
 \qquad
 \Delta_3=\alpha\gamma(3-x),
 \qquad
 \Delta_4=\alpha\gamma(4-x).
\end{equation}

\subsection{Price ranking in the baseline region $x<1$}

Using the closed forms,
\begin{equation}
 p^{IN}=\frac{\bar Q}{\alpha(2-x)},
 \qquad
 p^{VN}=\frac{2\bar Q}{\alpha(3-x)},
 \qquad
 p^{MS}=\frac{\bar Q(3-x)}{2\alpha(2-x)},
 \qquad
 p^{RS}=\frac{3\bar Q}{\alpha(4-x)}.
\end{equation}
To show $p^{RS}>p^{MS}$ when $x<1$, compare denominators:
\begin{align}
 p^{RS}>p^{MS}
 &\iff \frac{3}{4-x}>\frac{3-x}{2(2-x)} \\
 &\iff 6(2-x)>(3-x)(4-x) \\
 &\iff x(1-x)>0,
\end{align}
which holds for $0<x<1$.
Similarly,
\begin{align}
 p^{MS}>p^{VN}
 &\iff \frac{3-x}{2(2-x)}>\frac{2}{3-x} \\
 &\iff (3-x)^2>4(2-x) \\
 &\iff (x-1)^2>0,
\end{align}
which holds for all $x\neq 1$.
Finally,
\begin{align}
 p^{VN}>p^{IN}
 &\iff \frac{2}{3-x}>\frac{1}{2-x}
 \iff 2(2-x)>3-x
 \iff 1>x.
\end{align}
Hence for $x<1$,
\begin{equation}
 p^{RS}>p^{MS}>p^{VN}>p^{IN}.
\end{equation}

\subsection{Price ranking in the high-green-pull region $1<x<2$}

The same inequalities flip where the final comparison depends on the sign of
$1-x$. Since $1-x<0$ in this region,
\begin{equation}
 p^{IN}>p^{VN},
 \qquad p^{VN}>p^{RS},
 \qquad p^{MS}>p^{VN},
\end{equation}
which yields
\begin{equation}
 p^{IN}>p^{MS}>p^{VN}>p^{RS}.
\end{equation}

\subsection{Recycling-maturity ranking}

Using the closed forms,
\begin{equation}
 s^{IN}=\frac{\bar Q\beta}{\alpha\gamma(2-x)},
 \quad
 s^{VN}=\frac{\bar Q\beta}{\alpha\gamma(3-x)},
 \quad
 s^{MS}=\frac{\bar Q\beta}{2\alpha\gamma(2-x)},
 \quad
 s^{RS}=\frac{\bar Q\beta}{\alpha\gamma(4-x)}.
\end{equation}
Since the numerators are common and positive, rankings follow from comparing the
corresponding denominators. In particular,
\begin{equation}
 s^{IN}>s^{VN}>s^{RS}
 \qquad\text{for all }0<x<2.
\end{equation}
To compare $s^{VN}$ and $s^{MS}$,
\begin{align}
 s^{VN}>s^{MS}
 &\iff \frac{1}{3-x}>\frac{1}{2(2-x)} \\
 &\iff 2(2-x)>3-x \\
 &\iff 1>x.
\end{align}
Therefore,
\begin{equation}
 s^{IN}>s^{VN}>s^{MS}>s^{RS}
 \quad\text{if }x<1,
\end{equation}
and
\begin{equation}
 s^{IN}>s^{MS}>s^{VN}>s^{RS}
 \quad\text{if }1<x<2.
\end{equation}
The same logic applies to demand because each demand expression has the same
ordering structure after factoring out the common term $\alpha\gamma\bar Q$.

\subsection{Retailer-profit threshold between MS and RS}

The chapter states that in the high-green-pull region the ranking between
$\Pi_R^{MS}$ and $\Pi_R^{RS}$ flips at $x=4/3$. We verify that here.
Write
\begin{equation}
 \Pi_R^{MS}=\frac{\bar Q^2}{8\alpha(2-x)},
 \qquad
 \Pi_R^{RS}=\frac{\bar Q^2}{2\alpha(4-x)}.
\end{equation}
Then
\begin{align}
 \Pi_R^{MS}>\Pi_R^{RS}
 &\iff \frac{1}{8(2-x)}>\frac{1}{2(4-x)} \\
 &\iff 2(4-x)>8(2-x) \\
 &\iff 4-2x>16-8x \\
 &\iff 6x>12 \\
 &\iff x>\frac{4}{3}.
\end{align}
Hence
\begin{equation}
 \Pi_R^{RS}>\Pi_R^{MS}
 \quad\text{if }1<x<\frac{4}{3},
 \qquad
 \Pi_R^{MS}>\Pi_R^{RS}
 \quad\text{if }\frac{4}{3}<x<2.
\end{equation}

\section{Derivatives Used in the Comparative Statics Section}
\label{app:comparative_statics}

This section records the derivative calculations used in the comparative-statics
summary of Chapter~\ref{chap:model1}.

\subsection{Vertical Nash derivatives}

Using
\begin{equation}
 w^{VN}=m^{VN}=\frac{\bar Q\gamma}{3\alpha\gamma-\beta^2},
 \qquad
 s^{VN}=\frac{\bar Q\beta}{3\alpha\gamma-\beta^2},
\end{equation}
we obtain
\begin{equation}
 \frac{\partial w^{VN}}{\partial \alpha}
 =-\frac{3\bar Q\gamma^2}{(3\alpha\gamma-\beta^2)^2}<0,
 \qquad
 \frac{\partial s^{VN}}{\partial \alpha}
 =-\frac{3\bar Q\beta\gamma}{(3\alpha\gamma-\beta^2)^2}<0.
\end{equation}
Similarly,
\begin{equation}
 \frac{\partial w^{VN}}{\partial \beta}
 =\frac{2\bar Q\beta\gamma}{(3\alpha\gamma-\beta^2)^2}>0,
 \qquad
 \frac{\partial s^{VN}}{\partial \beta}
 =\frac{\bar Q(3\alpha\gamma+\beta^2)}{(3\alpha\gamma-\beta^2)^2}>0.
\end{equation}
For retailer profit,
\begin{equation}
 \Pi_R^{VN}=\frac{\gamma\bar Q^2(2\alpha\gamma-\beta^2)}{2(3\alpha\gamma-\beta^2)^2}.
\end{equation}
Differentiating with respect to $\alpha$ yields
\begin{equation}
 \frac{\partial \Pi_R^{VN}}{\partial \alpha}
 =-\frac{\bar Q^2\gamma^2(3\alpha\gamma-2\beta^2)}{(3\alpha\gamma-\beta^2)^3},
\end{equation}
which changes sign when $3\alpha\gamma-2\beta^2=0$, that is,
\begin{equation}
 \beta^2=\frac{3}{2}\alpha\gamma.
\end{equation}
Differentiating with respect to $\beta$ gives
\begin{equation}
 \frac{\partial \Pi_R^{VN}}{\partial \beta}
 =\frac{\bar Q^2\beta\gamma(\alpha\gamma-\beta^2)}{(3\alpha\gamma-\beta^2)^3},
\end{equation}
which changes sign at $\beta^2=\alpha\gamma$.
Finally,
\begin{equation}
 \frac{\partial \Pi_R^{VN}}{\partial \gamma}
 =-\frac{\bar Q^2\beta^2(\alpha\gamma-\beta^2)}{2(3\alpha\gamma-\beta^2)^3},
\end{equation}
which again changes sign at $\beta^2=\alpha\gamma$.

\subsection{Manufacturer-Stackelberg derivatives}

Because
\begin{equation}
 w^{MS}=\frac{\bar Q}{2\alpha},
\end{equation}
we immediately have
\begin{equation}
 \frac{\partial w^{MS}}{\partial \alpha}=-\frac{\bar Q}{2\alpha^2}<0,
 \qquad
 \frac{\partial w^{MS}}{\partial \beta}=0,
 \qquad
 \frac{\partial w^{MS}}{\partial \gamma}=0.
\end{equation}
The remaining closed forms in Chapter~\ref{chap:model1} imply the sign pattern
reported in the main text:
\begin{equation}
 \frac{\partial m^{MS}}{\partial \alpha}<0,
 \quad
 \frac{\partial s^{MS}}{\partial \alpha}<0,
 \quad
 \frac{\partial Q^{MS}}{\partial \alpha}<0,
 \quad
 \frac{\partial \Pi_M^{MS}}{\partial \alpha}<0,
 \quad
 \frac{\partial \Pi_R^{MS}}{\partial \alpha}<0,
\end{equation}
and
\begin{equation}
 \frac{\partial m^{MS}}{\partial \beta}>0,
 \quad
 \frac{\partial s^{MS}}{\partial \beta}>0,
 \quad
 \frac{\partial Q^{MS}}{\partial \beta}>0,
 \quad
 \frac{\partial \Pi_M^{MS}}{\partial \beta}>0,
 \quad
 \frac{\partial \Pi_R^{MS}}{\partial \beta}>0.
\end{equation}
Likewise,
\begin{equation}
 \frac{\partial m^{MS}}{\partial \gamma}<0,
 \quad
 \frac{\partial s^{MS}}{\partial \gamma}<0,
 \quad
 \frac{\partial Q^{MS}}{\partial \gamma}<0,
 \quad
 \frac{\partial \Pi_M^{MS}}{\partial \gamma}<0,
 \quad
 \frac{\partial \Pi_R^{MS}}{\partial \gamma}<0.
\end{equation}

\subsection{Retailer-Stackelberg derivatives}

For RS, the closed forms are
\begin{equation}
 w^{RS}=\frac{\gamma\bar Q}{4\alpha\gamma-\beta^2},
 \quad
 m^{RS}=\frac{2\gamma\bar Q}{4\alpha\gamma-\beta^2},
 \quad
 s^{RS}=\frac{\beta\bar Q}{4\alpha\gamma-\beta^2}.
\end{equation}
Therefore,
\begin{equation}
 \frac{\partial w^{RS}}{\partial \alpha}
 =-\frac{4\bar Q\gamma^2}{(4\alpha\gamma-\beta^2)^2}<0,
 \qquad
 \frac{\partial w^{RS}}{\partial \beta}
 =\frac{2\bar Q\beta\gamma}{(4\alpha\gamma-\beta^2)^2}>0,
\end{equation}
and
\begin{equation}
 \frac{\partial w^{RS}}{\partial \gamma}
 =-\frac{\bar Q\beta^2}{(4\alpha\gamma-\beta^2)^2}<0.
\end{equation}
The same denominator structure yields the sign pattern reported in the chapter:
all operational variables fall with $\alpha$, rise with $\beta$, and the
retailer and manufacturer profits fall with $\gamma$.

\section{Appendix to the Stage-0 Meta-Game}
\label{app:meta_game}

The chapter maps the stage-0 actions into operating regimes as follows:
\begin{equation}
 (C,C)\mapsto IN,
 \qquad
 (L,C)\mapsto MS,
 \qquad
 (C,L)\mapsto RS,
 \qquad
 (L,L)\mapsto VN.
\end{equation}
The argument for the baseline coordination game can be stated formally.

\subsection{Manufacturer ranking}

With the small non-cooperative overhead $F>0$, the manufacturer's net payoffs are
\begin{equation}
 \Pi_M^{VN}(F)=\Pi_M^{VN}(0)-F,
 \quad
 \Pi_M^{MS}(F)=\Pi_M^{MS}(0)-F,
 \quad
 \Pi_M^{RS}(F)=\Pi_M^{RS}(0)-F,
\end{equation}
and
\begin{equation}
 \Pi_M^{IN}(F)=\Pi_M^{IN}(0)-\frac{F}{2}.
\end{equation}
Since the chapter uses the equal-split normalization
$\Pi_M^{IN}(0)=\Pi_M^{MS}(0)$, it follows immediately that
\begin{equation}
 \Pi_M^{IN}(F)-\Pi_M^{MS}(F)=\frac{F}{2}>0.
\end{equation}
Therefore the manufacturer strictly prefers coordination to manufacturer
leadership once the tie-break term is introduced.

\subsection{Retailer ranking in the baseline region}

For $x<1$, the chapter's profit comparisons imply
\begin{equation}
 \Pi_R^{IN}>\Pi_R^{RS}>\Pi_R^{VN}>\Pi_R^{MS}.
\end{equation}
Hence the ordinal stage-0 payoff matrix is
\begin{center}
\begin{tabular}{c|cc}
 & $R:C$ & $R:L$ \\\hline
$M:C$ & $(4,4)$ & $(1,3)$ \\
$M:L$ & $(3,1)$ & $(2,2)$
\end{tabular}
\end{center}
The best responses are immediate:
\begin{itemize}
 \item if the retailer chooses $C$, then $M$ prefers $C$ because $4>3$;
 \item if the retailer chooses $L$, then $M$ prefers $L$ because $2>1$;
 \item if the manufacturer chooses $C$, then $R$ prefers $C$ because $4>3$;
 \item if the manufacturer chooses $L$, then $R$ prefers $L$ because $2>1$.
\end{itemize}
Thus $(C,C)$ and $(L,L)$ are Nash equilibria. The former is Pareto-superior,
while the latter is the inefficient coordination-failure outcome.

\subsection{High-green-pull region}

When $1<x<2$, the retailer's payoff from VN becomes relatively less attractive,
and the chapter shows that the retailer's ranking strengthens in favor of
coordination or, failing that, manufacturer leadership. This is why the set of
stage-0 equilibria narrows toward the coordinated outcome and the
manufacturer-led outcome with a willing retailer.
